\documentclass[11pt,letterpaper]{article}
\usepackage[margin=1in]{geometry}
\usepackage{amsmath, amssymb}
\usepackage{booktabs}
\usepackage{graphicx}
\usepackage{hyperref}
\usepackage{xcolor}
\usepackage{float}
\usepackage{array}
\usepackage{multirow}
\usepackage{caption}
\usepackage{subcaption}
\usepackage{parskip}

\hypersetup{colorlinks=true, linkcolor=blue, urlcolor=blue, citecolor=blue}

\title{\textbf{CSC311 Machine Learning Challenge Report}\\[0.4em]
       \large Painting Classification from Viewer Responses}
\author{%
  % TODO: Add all group member names and student numbers
  Group Member 1 \and Group Member 2 \and Group Member 3
}
\date{March 30, 2026}

\begin{document}
\maketitle
\tableofcontents
\newpage

% ─────────────────────────────────────────────────────────────────────────────
\section{Data Exploration and Feature Representation}
% ─────────────────────────────────────────────────────────────────────────────

\subsection{Dataset Overview}

The dataset contains 1{,}583 survey responses collected from viewers of three paintings:
\emph{The Persistence of Memory} (Dalí), \emph{The Starry Night} (van Gogh), and
\emph{The Water Lily Pond} (Monet). The task is multi-class classification: given a
viewer's responses, predict which painting they viewed. The class distribution is
roughly balanced across the three targets (approximately 500--550 responses per
painting), meaning accuracy is a natural and appropriate evaluation metric——a
random baseline would achieve $\approx 33\%$.

\subsection{Feature Inventory}

The raw CSV contains 15 columns (excluding \texttt{id} and \texttt{painting}):

\begin{table}[H]
\centering
\caption{Feature types in the dataset.}
\begin{tabular}{lll}
\toprule
\textbf{Feature} & \textbf{Type} & \textbf{Description}\\
\midrule
\texttt{emotion\_intensity} & Continuous  & Self-rated intensity of emotion (0--10)\\
\texttt{feel\_sombre}       & Ordinal int & How sombre the painting felt\\
\texttt{feel\_content}      & Ordinal int & How content the painting felt\\
\texttt{feel\_calm}         & Ordinal int & How calm the painting felt\\
\texttt{feel\_uneasy}       & Ordinal int & How uneasy the painting felt\\
\texttt{prominent\_colours} & Discrete int& Number of prominent colours noticed\\
\texttt{objects\_noticed}   & Discrete int& Number of distinct objects noticed\\
\texttt{willingness\_to\_pay}& Continuous & Amount willing to pay for the painting\\
\midrule
\texttt{room}       & Multi-label categorical & Room(s) they would hang it in\\
\texttt{season}     & Multi-label categorical & Season(s) they associate with it\\
\texttt{view\_with} & Multi-label categorical & Who they would view it with\\
\midrule
\texttt{feeling\_description} & Free text & Open-ended description of feelings\\
\texttt{food\_association}    & Free text & Food they associate with the painting\\
\texttt{soundtrack}           & Free text & Soundtrack they would pair it with\\
\bottomrule
\end{tabular}
\end{table}

\noindent The \texttt{soundtrack} column contained 4 missing values; all missing
text fields were imputed with empty strings before vectorisation.

\subsection{Exploratory Data Analysis}

We explored the data using histograms of each numerical feature grouped by
painting (implemented in \texttt{data\_description.py}). Key observations:

\begin{itemize}
  \item \textbf{feel\_sombre}: \emph{Persistence of Memory} responses were
        noticeably more sombre on average, while \emph{Water Lily Pond} responses
        clustered at the lower end. This feature provided strong discriminative signal.
  \item \textbf{feel\_calm} and \textbf{feel\_content}: \emph{Water Lily Pond}
        viewers reported higher calmness and contentment, consistent with Monet's
        impressionist style.
  \item \textbf{feel\_uneasy}: \emph{Persistence of Memory} produced the highest
        unease scores, reflecting Dalí's surrealist imagery.
  \item \textbf{emotion\_intensity}: Broadly distributed across all paintings, with
        moderate discriminative power on its own.
  \item \textbf{willingness\_to\_pay}: Showed some variation across paintings but
        with high within-class variance, making it a weaker standalone predictor.
  \item \textbf{prominent\_colours} and \textbf{objects\_noticed}: Subtle
        differences exist; \emph{Starry Night} tended to produce higher colour counts.
  \item \textbf{room}, \textbf{season}, \textbf{view\_with}: Multi-label responses
        that still correlated with painting——e.g., \emph{Water Lily Pond} was more
        frequently placed in living rooms or viewed with family, while Dalí's painting
        was often selected for a study or viewed alone.
\end{itemize}

% TODO: Insert histogram figures here, e.g.:
% \begin{figure}[H]
%   \centering
%   \includegraphics[width=0.8\textwidth]{figures/feel_sombre_hist.png}
%   \caption{Distribution of \texttt{feel\_sombre} per painting.}
% \end{figure}

\subsection{Data Representation}

We constructed two distinct feature sets, motivated by the heterogeneous nature of
the columns:

\paragraph{Structured features.}
\begin{itemize}
  \item \emph{Numerical}: the 8 ordinal/continuous columns above. We standardised
        these (zero mean, unit variance) using statistics computed from the training
        split, to prevent scale from dominating distance- or gradient-based models.
  \item \emph{Categorical (multi-label)}: \texttt{room}, \texttt{season}, and
        \texttt{view\_with} each allow multiple answers separated by commas.  We
        applied \textbf{multi-hot binarisation} (similar to \texttt{MultiLabelBinarizer}):
        each unique tag becomes a binary column. For example, \texttt{room} produced
        columns for \emph{Bedroom}, \emph{Kitchen}, \emph{Living room}, etc.
\end{itemize}

\paragraph{Text features.}
The three free-text columns (\texttt{feeling\_description}, \texttt{food\_association},
\texttt{soundtrack}) were vectorised into numerical features. During initial model
exploration (\texttt{ensemble/ensemble.py}), we used a simple \textbf{Bag-of-Words}
(BoW) representation via \texttt{TfidfVectorizer} with default settings. After
selecting the model architecture, we upgraded to a refined \textbf{TF-IDF}
configuration for the hyperparameter sweep and final model:
\begin{itemize}
  \item \texttt{max\_features=100} per column (top 100 terms / bigrams by document
        frequency), giving 300 total text features.
  \item \texttt{sublinear\_tf=True}: replaces raw term count $c$ with $1+\log c$,
        dampening the influence of very frequent terms.
  \item \texttt{ngram\_range=(1,2)}: includes both unigrams and bigrams, capturing
        short phrases such as ``calm feeling'' or ``dark sky''.
  \item \texttt{strip\_accents="unicode"} and default L2 normalisation.
\end{itemize}

\subsection{Data Splitting}

All experiments used a fixed \textbf{70\,/\,15\,/\,15} split
(\texttt{random\_state=42}):

\begin{itemize}
  \item \textbf{Training set} (70\%, $n\approx1108$): used to fit model parameters.
  \item \textbf{Validation set} (15\%, $n\approx237$): used to compare models and
        tune hyperparameters without contaminating the test set.
  \item \textbf{Test set} (15\%, $n\approx238$): held out entirely until final
        evaluation. We report test accuracy only for our final chosen configurations,
        not for every hyperparameter trial.
\end{itemize}

The same fixed split was used across \emph{all} model comparisons, ensuring the
evaluation metrics are directly comparable.

% ─────────────────────────────────────────────────────────────────────────────
\section{Models Evaluated}
% ─────────────────────────────────────────────────────────────────────────────

We evaluated three families of models. In all cases, both feature variants
(structured only, and structured + BoW text) were tried. This initial exploration
used BoW text features and was implemented in \texttt{ensemble/ensemble.py}.

\subsection{Logistic Regression}
We used $\ell_2$-regularised multinomial logistic regression with $C = 10$
(inverse regularisation strength) and no intercept term.  Both feature sets
were standardised as described in Section~1.  This serves as a strong linear
baseline and is well-suited to sparse TF-IDF features.

\subsection{Multi-Layer Perceptron (MLP)}
We trained a two-hidden-layer MLP with architecture $(150, 50)$, using SGD with
a fixed learning rate of 0.1 and up to 1{,}000 iterations.  The same
normalisation pipeline as logistic regression was applied.  The MLP is a
non-linear model that can learn feature interactions, making it a reasonable
second baseline.

\subsection{Decision Tree Family}
We explored six tree-based models:
\begin{itemize}
  \item \textbf{Decision Tree} (depth~5 and unlimited): simple interpretable baseline.
  \item \textbf{Random Forest} (200 trees): bagging ensemble; reduces variance by
        averaging many uncorrelated trees.
  \item \textbf{Extra Trees} (200 trees): similar to Random Forest but with
        additional randomisation of split thresholds.
  \item \textbf{AdaBoost} (200 estimators, learning rate 0.5): boosting algorithm
        that sequentially up-weights misclassified examples.
  \item \textbf{Gradient Boosting} (200 estimators, learning rate 0.1, depth~3):
        sequential boosting that fits each tree to the residual errors of the
        previous ensemble, using gradients of the loss function.
\end{itemize}

GradientBoosting requires no explicit normalisation (trees are invariant to
monotone feature rescaling), so raw structured features (before standardisation)
and raw TF-IDF vectors were used for all tree methods.

\subsection{Feature-Partitioned Ensemble}
After observing that the structured and text features carry complementary
information, we trained \emph{specialised} models on each feature set separately
and combined their predicted class probabilities via \textbf{soft voting}
(weighted average of probability vectors from both branches).
This was done for every combination of the 7~model types above ($7\times7 = 49$
combinations) using equal weights (\texttt{ensemble/ensemble.py}).

% ─────────────────────────────────────────────────────────────────────────────
\section{Model Selection and Hyperparameters}
% ─────────────────────────────────────────────────────────────────────────────

\subsection{Evaluation Protocol}

We use \textbf{accuracy} (fraction of correct predictions) as our primary metric.
Accuracy is appropriate because the three classes are approximately balanced
($\approx500$--$550$ examples each), so it is not misleadingly inflated by
majority-class prediction.

All model comparisons use \emph{the same} fixed \textbf{70/15/15}
train/validation/test split (\texttt{random\_state=42}), ensuring that accuracy
figures are directly comparable across experiments. Hyperparameters were selected
on \emph{validation} accuracy only; the test set was consulted only for the
final model selection table to estimate generalisation.

\subsection{Single-Model Results}

Table~\ref{tab:single} reports validation and test accuracy for all models under
both feature sets. The ``+ BoW'' column indicates whether BoW text features
were concatenated with the structured features. These experiments were conducted
with \texttt{ensemble/ensemble.py}.

\begin{table}[H]
\centering
\caption{Single-model accuracy (BoW text features). All experiments used the same
train/valid/test split.}
\label{tab:single}
\begin{tabular}{llrr}
\toprule
\textbf{Model} & \textbf{Features} & \textbf{Valid Acc} & \textbf{Test Acc}\\
\midrule
Decision Tree (depth=5)   & Structured only & 0.8186 & 0.8277 \\
Decision Tree (depth=5)   & + BoW            & 0.7764 & 0.8025 \\
Decision Tree (unlimited) & Structured only & 0.8017 & 0.7941 \\
Decision Tree (unlimited) & + BoW            & 0.7637 & 0.7983 \\
\midrule
Random Forest             & Structured only & 0.8734 & 0.8824 \\
Random Forest             & + BoW            & 0.8650 & 0.8824 \\
Extra Trees               & Structured only & 0.8565 & 0.8697 \\
Extra Trees               & + BoW            & 0.8734 & 0.8824 \\
AdaBoost                  & Structured only & 0.8861 & 0.8866 \\
AdaBoost                  & + BoW            & 0.8987 & 0.8950 \\
\textbf{Gradient Boosting}& \textbf{Structured only} & \textbf{0.8945} & \textbf{0.9118}\\
Gradient Boosting         & + BoW            & 0.9156 & 0.8950 \\
\midrule
Logistic Regression       & Structured only & 0.8565 & 0.8782 \\
Logistic Regression       & + BoW            & 0.8650 & 0.8908 \\
\midrule
MLP (150, 50)             & Structured only & 0.8819 & 0.8866 \\
MLP (150, 50)             & + BoW            & 0.8819 & 0.8613 \\
\bottomrule
\end{tabular}
\end{table}

\noindent Key observations:
\begin{itemize}
  \item \textbf{Gradient Boosting} achieves the best single-model test accuracy
        (91.2\%) on structured features alone.
  \item Concatenating BoW features into a single feature vector helps Logistic
        Regression (+1.3\%) and AdaBoost (+0.8\%), but \emph{hurts} MLP ($-2.5\%$)
        and does not reliably help GradientBoosting on the test set, suggesting the
        high-dimensional sparse text features can cause overfitting when naively
        concatenated.
  \item All tree-ensemble methods significantly outperform single decision trees,
        confirming the benefit of ensembling.
\end{itemize}

\subsection{Feature-Partitioned Ensemble Results}

Because simply concatenating BoW features into the structured matrix often
hurt or failed to help, we instead trained \emph{separate} specialised models
on each feature group and combined their class probabilities via soft voting.
Table~\ref{tab:ens} shows the top combinations (from \texttt{ensemble/ensemble.py},
using BoW text features).

\begin{table}[H]
\centering
\caption{Top-10 feature-partitioned ensemble combinations (equal-weight soft voting,
BoW text features).}
\label{tab:ens}
\begin{tabular}{llrr}
\toprule
\textbf{Structured model} & \textbf{Text model} & \textbf{Valid Acc} & \textbf{Test Acc}\\
\midrule
\textbf{Gradient Boosting} & \textbf{Gradient Boosting} & \textbf{0.9283} & \textbf{0.9412}\\
Gradient Boosting & Random Forest     & 0.9241 & 0.9202 \\
Logistic Regression & Gradient Boosting & 0.8945 & 0.9202 \\
Gradient Boosting & Logistic Regression & 0.8861 & 0.9160 \\
Gradient Boosting & Extra Trees       & 0.9156 & 0.9160 \\
MLP               & Logistic Regression & 0.9030 & 0.9160 \\
MLP               & Random Forest     & 0.9030 & 0.9118 \\
AdaBoost          & AdaBoost          & 0.8987 & 0.9118 \\
Gradient Boosting & AdaBoost          & 0.9114 & 0.9118 \\
MLP               & Gradient Boosting & 0.9156 & 0.9118 \\
\bottomrule
\end{tabular}
\end{table}

The GradientBoosting + GradientBoosting ensemble achieved the highest accuracy on
both validation (92.8\%) and test (94.1\%), confirming that: (i)~splitting the
feature space lets each model specialise, and (ii)~GradientBoosting is the
strongest learner for both feature groups. Having settled on this
\textbf{architecture} (GBM~+~GBM), we then upgraded the text representation from
BoW to the refined TF-IDF configuration described in Section~1.4, and proceeded
to tune hyperparameters on the full ensemble.

\subsection{Ensemble Hyperparameter Sweep}

With the GBM~+~GBM architecture and the upgraded TF-IDF text features in place,
we performed a full grid search over four hyperparameters using
\texttt{hp\_sweep.py}, applying the same GBM settings to \emph{both} ensemble
branches simultaneously:
\begin{itemize}
  \item \texttt{n\_estimators} $\in \{50, 100, 200\}$
  \item \texttt{learning\_rate} $\in \{0.05, 0.1, 0.2\}$
  \item \texttt{max\_depth} $\in \{1, 3, 5\}$
  \item \texttt{w\_struct} $\in \{0.3, 0.5, 0.7\}$ — ensemble weight for the
        structured branch (text weight $= 1 - \text{w\_struct}$)
\end{itemize}
This yields $3 \times 3 \times 3 \times 3 = 81$ configurations.
All 81 were evaluated on the same 70/15/15 split; hyperparameters were chosen
by \emph{validation accuracy only}.

\paragraph{Effect of \texttt{n\_estimators} and \texttt{learning\_rate}.}
Table~\ref{tab:nest-lr} fixes \texttt{max\_depth}=5 and \texttt{w\_struct}=0.5
to isolate the interaction between the number of boosting rounds and the step size.

\begin{table}[H]
\centering
\caption{Validation accuracy by \texttt{n\_estimators} and \texttt{learning\_rate}
(\texttt{depth}=5, \texttt{w\_struct}=0.5).}
\label{tab:nest-lr}
\begin{tabular}{cccc}
\toprule
 & \texttt{lr}=0.05 & \texttt{lr}=0.1 & \texttt{lr}=0.2 \\
\midrule
\texttt{n\_est}=50  & 0.8903 & 0.9198 & 0.9114 \\
\texttt{n\_est}=100 & 0.9198 & \textbf{0.9325} & 0.8987 \\
\texttt{n\_est}=200 & 0.9241 & 0.9114 & 0.8945 \\
\bottomrule
\end{tabular}
\end{table}

\noindent A moderate learning rate (0.1) with 100 estimators gives the best
validation accuracy (93.3\%).  Increasing to 200 estimators at \texttt{lr}=0.1
begins to overfit (train accuracy reaches 100\% while validation drops).
A fast learning rate (0.2) similarly overfits with deeper trees.

\paragraph{Effect of \texttt{max\_depth}.}
Table~\ref{tab:depth} fixes \texttt{lr}=0.1 and \texttt{w\_struct}=0.5 to
show how tree depth affects the bias--variance trade-off.

\begin{table}[H]
\centering
\caption{Train and validation accuracy by \texttt{max\_depth}
(\texttt{lr}=0.1, \texttt{w\_struct}=0.5).}
\label{tab:depth}
\begin{tabular}{ccccc}
\toprule
\texttt{depth} & \texttt{n\_est}=50 & \texttt{n\_est}=100 & \texttt{n\_est}=200 \\
\midrule
1 & 0.8819 & 0.8987 & 0.9072 \\
3 & 0.9072 & 0.9241 & 0.9198 \\
5 & 0.9198 & \textbf{0.9325} & 0.9114 \\
\bottomrule
\end{tabular}
\end{table}

\noindent Depth 5 with 100 estimators is the sweet spot: it
provides enough capacity to fit the signal without underfitting (depth 1), and
\texttt{n\_est}=100 prevents the overfitting that occurs at \texttt{n\_est}=200.

\paragraph{Effect of ensemble weight.}
Table~\ref{tab:wstruct} fixes \texttt{n\_est}=100, \texttt{lr}=0.1,
\texttt{depth}=5 to compare weighting strategies.

\begin{table}[H]
\centering
\caption{Validation accuracy by \texttt{w\_struct}
(\texttt{n\_est}=100, \texttt{lr}=0.1, \texttt{depth}=5).}
\label{tab:wstruct}
\begin{tabular}{ccc}
\toprule
\texttt{w\_struct} & \textbf{Train Acc} & \textbf{Valid Acc} \\
\midrule
0.3 & 0.9991 & 0.8819 \\
\textbf{0.5} & \textbf{1.0000} & \textbf{0.9325} \\
0.7 & 1.0000 & 0.9030 \\
\bottomrule
\end{tabular}
\end{table}

\noindent Equal weighting (\texttt{w\_struct}=0.5) outperforms both alternatives,
suggesting that the structured and text branches carry roughly equal predictive
value and should be trusted equally.

\paragraph{Top-5 overall configurations.}
Table~\ref{tab:top5} shows the top results across all 81 configurations, ranked
by validation accuracy.

\begin{table}[H]
\centering
\caption{Top-5 configurations from the full 81-point grid search.}
\label{tab:top5}
\begin{tabular}{ccccrr}
\toprule
\texttt{n\_est} & \texttt{lr} & \texttt{depth} & \texttt{w\_struct} & \textbf{Train} & \textbf{Valid}\\
\midrule
\textbf{100} & \textbf{0.1} & \textbf{5} & \textbf{0.5} & \textbf{1.0000} & \textbf{0.9325}\\
200  & 0.2  & 3 & 0.5  & 1.0000 & 0.9283 \\
100  & 0.1  & 3 & 0.5  & 0.9838 & 0.9241 \\
200  & 0.05 & 3 & 0.5  & 0.9838 & 0.9241 \\
200  & 0.05 & 5 & 0.5  & 1.0000 & 0.9241 \\
200  & 0.1  & 3 & 0.7  & 0.9964 & 0.9241 \\
\bottomrule
\end{tabular}
\end{table}

\subsection{Final Model}

Based on the sweep results, our final submission uses a
\textbf{GradientBoosting + GradientBoosting weighted-average ensemble}:

\begin{itemize}
  \item \textbf{Structured branch}: GradientBoostingClassifier with
        \texttt{n\_estimators=100}, \texttt{learning\_rate=0.1}, \texttt{max\_depth=5},
        trained on the 8~standardised numerical columns plus multi-hot room, season,
        and view\_with features.
  \item \textbf{Text branch}: same GBM hyperparameters, trained on 300-dimensional
        TF-IDF features from \texttt{feeling\_description}, \texttt{food\_association},
        and \texttt{soundtrack}.
  \item \textbf{Combination}: equal-weight soft vote ($w_{\text{struct}} = 0.5$).
        The final probability vector is
        $p = 0.5\, p_{\text{struct}} + 0.5\, p_{\text{text}}$, and the predicted
        class is $\arg\max p$.
  \item Both models are trained on the \textbf{full available dataset} (all 1{,}583
        rows) at submission time, to maximise generalisation.
\end{itemize}

The parameters (tree node arrays, TF-IDF vocabulary and IDF weights, normalisation
statistics) are serialised to \texttt{model\_params.npz}. Inference in
\texttt{pred.py} uses only \texttt{numpy}, \texttt{pandas}, and Python standard
library modules; no \texttt{sklearn} is required at runtime.

% ─────────────────────────────────────────────────────────────────────────────
\section{Expected Prediction Performance}
% ─────────────────────────────────────────────────────────────────────────────

\paragraph{Point estimate.} We expect our model to achieve approximately
\textbf{91--93\% accuracy} on the unseen test set.

\paragraph{Reasoning.}
Our best ensemble configuration achieved 93.3\% validation accuracy on the
15\% held-out split. However, we note two factors that may push generalisation
performance slightly below that figure:

\begin{enumerate}
  \item \textbf{Training on full data.} The submitted model is retrained on all
        1{,}583 examples (no held-out set). Training on more data generally
        \emph{improves} generalisation, but makes it impossible to measure
        overfitting directly using our own split.
  \item \textbf{Distribution shift.} The competition test set was collected
        independently. Viewers of the same paintings may exhibit different response
        styles (e.g., different phrasings in the text fields). TF-IDF vocabulary
        coverage may be lower, moderately degrading the text branch.
\end{enumerate}

The structured numerical/categorical features are likely to transfer well (they are
scale-free and domain-agnostic), while the TF-IDF text branch may generalise
less perfectly. Given that the structured-only GradientBoosting achieved 91.2\%
on our test split and the ensemble consistently adds 2--3 percentage points of
complementary signal, we estimate \textbf{91--93\%} as a well-supported range.

% ─────────────────────────────────────────────────────────────────────────────
\section{Workload Distribution}
% ─────────────────────────────────────────────────────────────────────────────

% Each student must write their own contribution description.

\begin{itemize}
  \item \textbf{[Name, Student \#]}: % TODO: written by this student
  \item \textbf{[Name, Student \#]}: % TODO: written by this student
  \item \textbf{[Name, Student \#]}: % TODO: written by this student
\end{itemize}

\end{document}
